\section{Round-nine reconstruction: closed nullspaces, a scale connection, and resolved-memory covariance}
\label{sec:r9-c2}

Three objects are kept distinct: the spectral projector selecting a prepared
phase, the finite-dimensional resolved projection used in memory, and the
nonlinear observable map used in contraction.  Parameter transport is
constructed on a common strong--weak transfer scale and never uses
Radon--Nikodym densities between infinite-volume equilibrium measures, which
may be mutually singular.

\subsection{The two closed linear nullspaces}

Let $\mathscr A_W$ be the weighted dynamically H\"older potential algebra on
one labelled platform.  Define
\[
 \mathscr N_{\rm int}
 =\overline{\operatorname{span}}^{\mathscr A_W}
 \{U-U\circ\Theta_t:U\in\mathscr A_W,
   t\in\mathbb Q_+\},
\]
and
\[
 \mathscr N_{\rm pr}=\mathbb R\mathbf1+\mathscr N_{\rm int}.
\]
The first space is the nullspace for invariant integrals; the second is the
nullspace for normalized pressure tangents.

\begin{lemma}[Uniform weighted Ces\`aro tightness]
\label{lem:r9-c2-cesaro}
On each platform there is a proper Lyapunov weight $W$ such that for every
finite signed weighted Radon measure $\mu$,
\[
 \sup_{T\ge1}
 \frac1T\int_0^T
  \int W\circ\Theta_t\,d|\mu|\,dt
 \le C\int(1+W)d|\mu|.
\]
Consequently the Ces\`aro orbit averages of $\mu$ are relatively compact in
the weighted strict dual.
\end{lemma}

\begin{proof}
For Sinai suspensions use the exponential return/roof Lyapunov function from
A3/A4 and split at renewals; stationarity of the induced Gibbs kernel gives a
uniform average drift bound.  For hard spheres use conservation of mass and
energy plus the B2 exponential velocity/contact weight.  The A1 benchmark has
bounded state and is immediate.  Markov's inequality on the compact sublevels
of $W$ gives uniform tightness, while the displayed bound gives uniform
weighted integrability.
\end{proof}

\begin{theorem}[Invariant separator and Hausdorff cotangents]
\label{thm:r9-c2-null}
For $F\in\mathscr A_W$,
\[
 \nu(F)=0\quad\text{for every invariant finite-weight probability }\nu
 \quad\Longleftrightarrow\quad
 F\in\mathscr N_{\rm int}.
\]
Thus $\mathscr A_W/\mathscr N_{\rm int}$ and
$\mathscr A_W/\mathscr N_{\rm pr}$ are Hausdorff locally convex spaces with
the advertised invariant-integral and normalized-pressure meanings.
\end{theorem}

\begin{proof}
The forward implication for coboundaries is invariance and closure.  If
$F\notin\mathscr N_{\rm int}$, Hahn--Banach gives a finite signed weighted
Radon separator $\mu$ because the continuous dual of the weighted strict
space is the countably additive weighted Radon space.  Average $\mu$ over
$[0,T]$.  Lemma~\ref{lem:r9-c2-cesaro} gives a convergent subsequence; the
limit is invariant and keeps the separating pairing because the averaged
pairing of every rational-time coboundary vanishes.  Its Jordan parts and the
constant test produce invariant probabilities separating $F$.  Adding
constants gives the pressure quotient.
\end{proof}

\subsection{A Kato connection on a Banach scale}

Let
\[
 \mathbb B^{m+2}\hookrightarrow\mathbb B^{m+1}
 \hookrightarrow\mathbb B^{m}
\]
be the common moving-cut/trace or cluster transfer scale supplied by A2 or B2.
A phase chart has transfer/Feynman--Kac generator $L_\eta$ with
\[
 \partial_\eta L_\eta:
 \mathbb B^{m+1}\longrightarrow\mathbb B^{m},
\]
and a simple Riesz projector $\Pi_\eta^{\rm sp}$ on the strong-to-weak
resolvent contour.  Its derivative is the bounded scale map
\[
 (\Pi_\eta^{\rm sp})'
 =\frac1{2\pi i}\int_\Gamma
 (z-L_\eta)^{-1}L_\eta'(z-L_\eta)^{-1}dz.
\]

Define the scale connection on a section $u_\eta\in\operatorname{Ran}
\Pi_\eta^{\rm sp}$ by
\[
 \nabla_\eta^{\rm sp}u
 =\Pi_\eta^{\rm sp}\partial_\eta u
 +[\Pi_\eta^{\rm sp},(\Pi_\eta^{\rm sp})']u,
\]
where the derivative is interpreted one level weaker.  Equivalently, parallel
transport solves Kato's ODE
\[
 U_\eta=[(\Pi^{\rm sp})',\Pi^{\rm sp}]U,
 \qquad U(0)=I,
\]
on the scale.

\begin{theorem}[Scale Kato transport of the full phase data]
\label{thm:r9-c2-kato}
The Kato ODE has an invertible evolution from $\mathbb B^{m+1}$ to
$\mathbb B^m$ and transports the phase eigenline.  In the resulting common
trivialization, the full Doob generator
\[
 L_\eta^D=h_\eta^{-1}
 (L_\eta-P_\eta)h_\eta.
\]
The invariant left functional, the pairing, and every declared resolved
observable family are differentiable scale sections.  No density
$d\nu_\eta/d\nu_0$ on the infinite path space is used.
\end{theorem}

\begin{proof}
The resolvent formula and the one-level derivative loss make
$[(\Pi^{\rm sp})',\Pi^{\rm sp}]$ integrable as a bounded map between adjacent
levels.  Picard iteration on the finite parameter interval gives $U$ and its
inverse.  Analytic perturbation supplies $h_\eta$, the left
eigenfunctional, and $P_\eta$.  Conjugating each finite-time
transfer operator by $U$ gives differentiability of the full Doob family and
its pairings.  This construction remains meaningful even when distinct
stationary path measures are mutually singular.
\end{proof}

\subsection{Spectral and resolved projections are different}

Let $\Pi_\eta^{\rm sp}$ select the prepared stationary phase.  In
the phase Hilbert space obtained from its finite-time correlation pairing, let
$R_\eta$ be the orthogonal projection onto a declared finite family
of resolved observables.  Generally
$R_\eta\ne\Pi_\eta^{\rm sp}$.  The latter would project
onto constants and would indeed produce no nontrivial memory.

Transport the resolved basis by the Kato trivialization and re-orthonormalize
with its differentiable Gram matrix.  This defines $R_\eta$ and
$\nabla R_\eta$ without a measure-density trivialization.

\begin{theorem}[Covariant compressed resolvent and memory]
\label{thm:r9-c2-memory}
On the resolved range define
\[
 C_\eta(z)
 =R_\eta(z-L_\eta^D)^{-1}R_\eta,
 \qquad
 A_\eta=R_\eta L_\eta^D R_\eta,
\]
and
\[
 \widehat K_\eta(z)
 =zR_\eta-A_\eta-C_\eta(z)^{-1}.
\]
These are differentiable scale sections on the common right half-plane.  Their
covariant derivative is
\[
\begin{aligned}
 \nabla\widehat K
 ={}&z\nabla R-\nabla A
 +C^{-1}(\nabla C)C^{-1},\\
 \nabla C
 ={}&(\nabla R)(z-L^D)^{-1}R
 +R(z-L^D)^{-1}(\nabla R)\\
 &+R(z-L^D)^{-1}(\nabla L^D)(z-L^D)^{-1}R.
\end{aligned}
\]
The formula is the derivative of the actual resolved memory, not of the
spectral eigenprojector.
\end{theorem}

\begin{proof}
The A4 full-$Q$ compressed-resolvent theorem gives invertibility of $C(z)$ on
the declared half-plane.  Differentiate the resolvent identity in the Kato
trivialization and use
$\nabla(C^{-1})=-C^{-1}(\nabla C)C^{-1}$.  Substitution in
$zR-A-C^{-1}$ gives the displayed signs.  The derivative of the resolved Gram
projection is included in both outer terms.
\end{proof}

\subsection{Optional projections and nonlinear contraction}

Let $\mathcal F_t^{\rm coarse}$ be the genuine A4 past/coarse filtration.
For a terminal bounded functional $G$, the finite optional projection is
\[
 M_t^\varepsilon
 =\mathbb E_\varepsilon[G\mid\mathcal F_t^{\rm coarse}].
\]

\begin{theorem}[Stable optional-projection limit]
\label{thm:r9-c2-optional}
On a fixed-support phase chart, the A4 quenched kernel theorem implies joint
stable convergence of $M^\varepsilon$ with the enhanced resolved process
to the diffusion optional projection.  The conclusion is uniform for the
finite family of Kato-transported resolved observables used in
Theorem~\ref{thm:r9-c2-memory}.
\end{theorem}

\begin{proof}
Represent $M_t^\varepsilon$ as the integral of $G$ against the A4 past
kernel.  Uniform kernel convergence on compact histories gives convergence of
the integral and multiplication by bounded initial-history variables.
Exponential containment removes the complement.  Fixed support ensures all
finite-time laws have one common cylinder reference; no infinite-volume RN
derivative is asserted.
\end{proof}

Let $\mathcal C:X\to Y$ be a nonlinear continuous observable map between
labelled platform state spaces.  A cotangent test $\varphi$ on $Y$ pulls back
as $\varphi\circ\mathcal C$; at a differentiability point its tangent
pullback is $D\mathcal C(x)^*D\varphi(\mathcal Cx)$.

\begin{theorem}[Typed nonlinear contraction and tangent chain rule]
\label{thm:r9-c2-main}
If the parent laws have the good LDP proved on their labelled platform and
$\mathcal C$ is continuous or exponentially approximable, then
\[
 I_Y(y)=\inf_{\mathcal Cx=y}I_X(x).
\]
Pressure tests are composed with $\mathcal C$, and exposed tangent/covariance
maps obey the ordinary derivative and adjoint chain rules in the declared
source manifolds.  Combined with Theorems~\ref{thm:r9-c2-null}--
\ref{thm:r9-c2-optional}, this yields well-typed cotangent, memory, likelihood,
and diffusion diagrams without identifying distinct physical platforms.
\end{theorem}

\begin{proof}
The first statement is the extended contraction principle.  At finite volume,
$\langle\varphi,\mathcal C(X_\varepsilon)\rangle$ is exactly the parent
test $\varphi\circ\mathcal C$, so differentiation gives the chain rule when
the map is smooth.  Exponentially good approximations pass it to the limit.
The previous theorems supply the closed cotangent quotient, common scale
connection, resolved-memory derivative, and optional-projection limit on each
labelled component.
\end{proof}
